\renewcommand{\refname}{Referencias y recursos utilizados}
\begin{thebibliography}{7}
\setlength{\itemsep}{3pt}
\bibitem{tema11} MUCC--UPV. \emph{Enunciado del trabajo y Tema 11: Criptografía postcuántica}.
Criptología y Seguridad de los Datos, 2026. Material de la asignatura.
\bibitem{ggh} O.~Goldreich, S.~Goldwasser y S.~Halevi. Public-key cryptosystems from
lattice reduction problems. \emph{CRYPTO '97}, LNCS~1294, pp.~112--131, 1997.
\bibitem{galbraith} S.~D.~Galbraith. \emph{Mathematics of Public Key Cryptography}.
Cambridge University Press, 2012. Capítulos~17 y~18, versión ampliada y corregida:
\url{https://www.math.auckland.ac.nz/~sgal018/crypto-book/crypto-book.html}.
\bibitem{ia} OpenAI. \emph{Codex}. Página oficial, consulta: 16/11/2026.
\url{https://openai.com/codex/}.
\bibitem{claude} Anthropic. \emph{Claude Code}. Documentación oficial, consulta: 16/11/2026.
\url{https://code.claude.com/docs/en/overview}.
\bibitem{python} Python Software Foundation. \emph{Python 3: biblioteca estándar},
\texttt{fractions}, \texttt{math}, \texttt{time} y \texttt{statistics}.
\url{https://docs.python.org/3/library/}.
\bibitem{latex} LaTeX Project. \emph{LaTeX}; J.~Collins y colaboradores, \emph{latexmk};
C.~Feuersänger, \emph{PGFPlots}. Herramientas de composición de la memoria y su figura.
\end{thebibliography}

\noindent\textbf{Prompts y uso de IA.} Se emplearon OpenAI Codex y Claude Code
como apoyo técnico y de redacción.
Consultas representativas, reformuladas para facilitar su lectura:
\begin{enumerate}
\item Si la base no es ortogonal, ¿cómo buscamos el punto más cercano?
\item ¿Cómo hago el sistema en Python para sacar los multiplicadores?
\item ¿Para qué necesitamos Gram-Schmidt en estos algoritmos?
\item ¿Por qué al cambiar la base con LLL seguimos teniendo la misma celosía?
\item ¿Qué diferencia hay entre los dos métodos de Babai?
\item ¿Qué tengo que tener en cuenta para elegir el radio de búsqueda?
\item Si LLL y Babai no siempre dan el óptimo, ¿qué nos falta para asegurarlo?
\item ¿Puedes revisar el código por si se nos ha pasado algún caso?
\end{enumerate}
Las respuestas se usaron como propuestas de código, pruebas y redacción. Se contrastaron
casos pequeños y se ejecutaron las 19 tareas; las 250 pruebas actuales pasan.
Los resultados y tiempos proceden del programa. Biblioteca estándar y herramientas de
composición:~\cite{python,latex}.
